\documentclass[conference]{IEEEtran}
\IEEEoverridecommandlockouts
% The preceding line is only needed to identify funding in the first footnote. If that is unneeded, please comment it out.
\usepackage{cite}
\usepackage{amsmath,amssymb,amsfonts}
\usepackage{algorithmic}
\usepackage{graphicx}
\usepackage{textcomp}
\usepackage{xcolor}
\def\BibTeX{{\rm B\kern-.05em{\sc i\kern-.025em b}\kern-.08em
    T\kern-.1667em\lower.7ex\hbox{E}\kern-.125emX}}
\begin{document}

\title{AANEC-Relaxed Trapped Surface Theorem and Stellar Mass-Gap Derivation}\author{Anonymous Research Plan Author}\maketitle

\begin{abstract}
This research plan proposes a qualified AANEC-relaxed trapped-surface theorem for spherically symmetric stellar collapse with isotropic pressure, classical general relativity, global hyperbolicity, regular matter fields, complete null generators, and an explicit averaged null-convergence condition. The central contribution is a formal derivation chain that links Raychaudhuri-type focusing and Misner-Sharp compactness to trapped-surface formation, then uses a TOV-like stability predicate to exclude stable non-black-hole remnants above a model-dependent compactness threshold. The proposal explicitly separates this mathematical result from the observed 2-5 solar-mass gap; the gap is treated only as an astrophysical consistency constraint, not as a theorem-derived prediction, because explosion energetics, fallback, equation-of-state uncertainty, rotation, magnetic fields, nonspherical dynamics, pair-instability processes, and binary evolution can dominate remnant distributions. Quantum stress-energy violations of pointwise energy conditions motivate the averaged-condition framing, while the added focusing assumption remains an unverified geometric hypothesis rather than a consequence of AANEC alone.

The planned decision boundary is a proof or counterexample search. A closed proof would require the averaged-focusing and compactness obligations to hold for all admissible metrics in the restricted model class, yielding a proposed exclusion bound for stable non-black-hole remnants above the threshold. A valid falsifier would be a stable ultra-compact witness satisfying all declared assumptions, including the independent averaged-focusing condition, yet failing to form a trapped surface; proposed witnesses based on regular black-hole or metastable high-compactness geometries remain unverified because their satisfaction of the focusing and stability premises is not established. If the formal obligations cannot be closed, the theorem must be restricted or rejected within the declared domain, while any mass-gap inference remains non-discriminating when astrophysical alternatives dominate.
\end{abstract}
\begin{IEEEkeywords}
Averaged Null Energy Condition, trapped surfaces, stellar collapse, black hole remnants, compactness threshold, mass gap, general relativity
\end{IEEEkeywords}\section{Introduction}
\subsection{Motivation and Research Gap}
Classical general relativity traditionally relies on pointwise energy conditions to establish the inevitability of trapped-surface formation during gravitational collapse. However, quantum field theory demonstrates that pointwise conditions are systematically violated by realistic matter, necessitating a shift toward averaged energy conditions such as the Achronal Averaged Null Energy Condition (AANEC). While AANEC provides a robust lower bound for singularity theorems, it is currently an open research gap whether AANEC alone is sufficient to guarantee the local focusing required for trapped-surface formation without an explicit, independent averaged null-convergence assumption.~\cite{cite_0ed84243cf171bc4}

\subsection{Proposed Contribution}
This proposal introduces a qualified trapped-surface theorem for spherically symmetric stellar collapse. The central contribution is the derivation of a model-dependent compactness threshold that forces trapped-surface formation when AANEC is supplemented by an explicit averaged null-convergence condition. By establishing this geometric bound, the proposal aims to exclude stable non-black-hole remnants above a specific compactness scale within this restricted model class. This formal derivation treats the observed 2-5 solar mass gap not as a direct theorem prediction, but as an astrophysical consistency constraint subject to alternative stellar-physics explanations.

\subsection{Design Assumptions and Scope}
\begin{itemize}
\item The formal argument is bounded by the following design assumptions:
\item The spacetime is modeled as spherically symmetric with isotropic pressure, and matter fields satisfy AANEC along complete null generators.
\item An explicit averaged null-convergence condition is assumed as an independent geometric hypothesis, rather than a derived consequence of AANEC.
\item The precise mathematical forms of the compactness parameter and the averaged focusing condition are deferred to the definitions ledger; without them, the proof obligations cannot be closed.
\end{itemize}

\section{Background, Survey, and Research Gap}
The theoretical foundations of black hole formation rely heavily on singularity theorems that demand pointwise energy conditions. However, recent developments in quantum field theory indicate that these pointwise conditions are systematically violated by quantum fields. This necessitates the use of weaker, averaged statements to maintain broader validity. The AANEC has emerged as a candidate for a universal property of gravitating matter, offering a robust pathway to preserve singularity theorems even when local energy densities are negative. This shift from pointwise to averaged conditions is critical for reconciling classical general relativity with semiclassical effects.~\cite{cite_0ed84243cf171bc4}

Despite the promise of AANEC, a significant gap remains in the formal literature regarding the precise assumptions and validity domains under which Penrose-type singularity theorems and Einstein-field-equation solutions support the existence of black holes. Specifically, the required evidence slot for 'formal\_claim' is missing from the current survey evidence plan. Without explicitly stating the assumptions that bridge AANEC to trapped-surface formation, the claim cannot be fully evaluated for validity. This missing formal claim represents the primary unresolved research constraint.

In the astrophysical domain, the observed lower stellar mass gap between neutron stars and black holes provides a crucial constraint on formation mechanisms. Current models suggest this gap may arise from the maximum neutron-star mass, supernova explosion energy, fallback, rotation, magnetic fields, binary evolution, or observational selection effects. The existence of this gap serves as a discriminator for supernova engine physics, implying that the mass distribution is astrophysically inferred rather than a pure general-relativistic theorem.

A parallel gap exists in the identification of boundary variables that distinguish black-hole from non-black-hole outcomes. The required slot for 'boundary\_variable' remains uncovered in relevant evidence clusters concerning theoretical stellar evolution and remnant formation pathways. Identifying variables such as mass, compactness, rotation, or metallicity that serve as boundaries is essential for mapping the parameter space where black hole remnants are expected.

Furthermore, the derived mass distributions remain highly sensitive to uncertainties in explosion and unbinding energies across different progenitor masses. This leaves a key physical parameter unconstrained, introducing significant variance into the final remnant distribution. The interplay between these astrophysical uncertainties and the missing formal geometric assumptions highlights the need for a qualified AANEC trapped-surface bound that explicitly separates formal theorem constraints from astrophysical confounders.

\section{Research Questions and Planned Contributions}
\subsection{Core Research Questions}
This research program addresses the formal and astrophysical constraints governing stellar-mass black hole formation by proposing a qualified theorem that relaxes pointwise energy conditions. The central inquiry is whether an explicit averaged null-convergence condition, combined with the Achronal Averaged Null Energy Condition (AANEC), can rigorously force trapped-surface formation in spherically symmetric collapse. This inquiry is motivated by the observation that pointwise energy conditions are systematically violated by quantum fields, necessitating weaker statements such as averaged energy conditions for broader validity (W3009127927). The planned contribution establishes a model-dependent compactness threshold that excludes stable non-black-hole remnants within this restricted geometric class. Furthermore, the research questions delineate the boundary between formal theorem validity and astrophysical application, explicitly treating the observed 2-5 solar mass gap as an astrophysical consistency constraint rather than a direct theorem prediction, given that explosion energetics and fallback dynamics may independently govern remnant distributions.~\cite{cite_7a407c865faafd11}

\begin{table*}[htbp]
\centering
\footnotesize
\setlength{\tabcolsep}{3pt}
\renewcommand{\arraystretch}{1.12}
\caption{Planned Contributions and Decision Criteria}
\label{tab:research-questions-and-contributions-rq-contributions-table}
\begin{tabular}{|p{0.29\textwidth}|p{0.29\textwidth}|p{0.29\textwidth}|}
\hline
Contribution Target & Decision Criterion & Required Formal Dependency \\\hline
AANEC-Relaxed Trapped Surface Theorem & Closure of proof obligations PO1 and PO2 for all admissible metrics & Mathematical form of explicit averaged null-convergence condition (A5) \\\hline
Stable Remnant Exclusion Bound & Derivation of S3-S4 from stable equilibrium limit (A6) and compactness parameter (D4) & Formal definition of Misner-Sharp compactness C \\\hline
Counterexample Validity Audit & Verification that candidate witnesses CE1 and CE2 satisfy all assumptions A1-A6 & Human review of A5 definition and A6 stability predicate \\\hline
Mass-Gap Astrophysical Relevance & Demonstration that explosion energetics do not fully dominate observed remnant distribution & Separation of formal theorem from stellar-evolution confounders \\\hline
\end{tabular}
\end{table*}

\subsection{Scope and Unresolved Dependencies}
The intellectual labor of this section centers on establishing the precise scope of the proposed theorem while identifying the formal gaps that currently delimit its operationalization. The argument transitions to the definitions and propositions stage by exposing the critical dependency on the mathematical form of the explicit averaged null-convergence condition and the Misner-Sharp compactness parameter. Without these definitions, the stability of non-black-hole remnants cannot be formally distinguished from regular black hole models, and the status of candidate counterexamples remains ambiguous. The research design maintains a strict separation between formal and empirical claims: while the theorem addresses geometric focusing and compactness bounds, the astrophysical mass gap remains subject to alternative explanations such as rotation, magnetic fields, and binary evolution history. This distinction ensures that the planned contributions remain grounded in rigorous geometric reasoning rather than overreaching observational inference.

\section{Idea Source Checkpoints and Direction Selection Audit}
\subsection{Scientific attraction of the qualified AANEC route}
The selected direction attracts because it replaces fragile pointwise energy assumptions with an averaged condition while keeping the astrophysical mass gap outside the theorem’s scope. The proposal is that spherically symmetric collapse, if it satisfies AANEC, global hyperbolicity, regular matter fields, complete null generators, and an explicit averaged null-convergence/focusing condition, reaches a model-dependent compactness threshold that forces a trapped surface. Stable non-black-hole remnants above the corresponding compactness-dependent scale would then be excluded within that restricted model class. This framing treats the 2-5 solar mass gap as a consistency constraint rather than a theorem-derived prediction, preserving the formal claim from confounding by explosion energetics, fallback, equation-of-state uncertainty, rotation, magnetic fields, nonspherical dynamics, and binary evolution.

\begin{table*}[htbp]
\centering
\footnotesize
\setlength{\tabcolsep}{3pt}
\renewcommand{\arraystretch}{1.12}
\caption{Audit snapshot comparison}
\label{tab:idea-origin-and-selection-checkpoint-audit-table}
\begin{tabular}{|p{0.23\textwidth}|p{0.23\textwidth}|p{0.23\textwidth}|p{0.23\textwidth}|}
\hline
Checkpoint & Retained element & Exposed defect or correction & Route decision \\\hline
Candidate snapshot & AANEC plus averaged focusing; compactness threshold; mass gap demoted to astrophysical consistency & Misner-Sharp compactness and averaged focusing are not yet formalized & Keep as proposed mechanism, not proved theorem \\\hline
Portfolio snapshot & Qualified title and restricted non-black-hole-remnant exclusion & Stability predicate for exotic remnants remains undefined & Carry forward only under explicit A6-style stability assumption \\\hline
Direction snapshot & Independent averaged focusing condition; counterexample boundary; confounder list & Candidate witnesses cannot be adjudicated without A5 and C definitions & Advance to formal definition and proof-obligation ledger \\\hline
\end{tabular}
\end{table*}

\subsection{Defects that convert attraction into a qualified proposal}
The supplied checkpoints expose a decisive defect: AANEC alone is not treated as sufficient for the desired trapped-surface conclusion. Survey evidence on energy conditions supports the motivation for averaged statements because pointwise conditions can fail in quantum regimes, but the route still requires a separate averaged focusing hypothesis to close the Raychaudhuri-to-trapped-surface inference. The candidate counterexamples illustrate the same problem. A high-compactness regular object or a regular black-hole model may evade the exclusion only if it fails the added focusing condition, completeness, regularity, or the stability predicate. Therefore, the retained direction must be qualified as a conjecture with explicit validity domains. The next section should define the compactness symbol, state the averaged focusing condition, and separate formal proof obligations from astrophysical explanations.~\cite{cite_0ed84243cf171bc4}

\section{Problem Definition, Assumptions, and Hypotheses}
\subsection{Problem Definition and Formal Scope}
The formal problem is to determine whether an averaged null-energy hypothesis can replace pointwise energy conditions in a trapped-surface theorem for spherically symmetric stellar collapse, and whether the resulting compactness bound excludes stable non-black-hole remnants. The admissible domain is restricted to classical general relativity with regular matter fields, spherical symmetry, isotropic pressure, global hyperbolicity, complete null generators, and a proposed averaged focusing condition. The problem is not a quantitative derivation of the observed stellar mass gap, nor a general proof that every massive star forms a black hole; astrophysical mass-gap explanations remain separate empirical claims governed by explosion energetics, fallback, equation-of-state physics, rotation, magnetic fields, nonspherical dynamics, and binary evolution.

\subsection{AANEC and Averaged Focusing Relation}
∫\_γ T\_ab k\textasciicircum{}a k\textasciicircum{}b dλ ≥ 0~\cite{cite_0ed84243cf171bc4}

\subsection{Meaning of the Proposed Relation}
The displayed relation is the proposed AANEC-style constraint along a null generator γ with tangent k\textasciicircum{}a, stress-energy tensor T\_ab, and affine parameter λ. It is intended to fix the matter assumption consumed by the later Raychaudhuri/focusing chain. The relation alone is not sufficient for the theorem as proposed: the plan additionally requires an explicit averaged null-convergence or focusing condition on the relevant congruence. That extra condition is the central unresolved formal dependency and must not be collapsed into the AANEC integral.~\cite{cite_0ed84243cf171bc4}

\subsection{Assumption Ledger for This Route}
\begin{itemize}
\item A1: Classical general relativity, spherical collapse, and isotropic pressure.
\item A2: Matter satisfies AANEC along complete null generators.
\item A3: Global hyperbolicity and sufficient regularity for trapped-surface arguments.
\item A4: Completeness of the relevant null generators.
\item A5: An explicit averaged null-convergence or focusing condition holds along the relevant generators.
\item A6: Stable non-black-hole remnants obey a bounded compactness or TOV-like equilibrium limit.
\end{itemize}

\subsection{Proposed Problem Proposition}
\paragraph{Proposition (proposed).} Proposition P1 (proposed): In the restricted model class defined by A1-A5, if a spherically symmetric collapse reaches a model-dependent compactness threshold, then a trapped surface forms. Proof obligation: close the Raychaudhuri focusing step and the trapped-surface identification step for all admissible metrics in the restricted class. Failure condition: the proposition fails or must be restricted if AANEC is violated, the added averaged focusing condition fails, null generators are incomplete, quantum stress-energy effects prevent the needed focusing, or nonspherical, rotating, magnetic, or anisotropic dynamics dominate the collapse.

\begin{table*}[htbp]
\centering
\footnotesize
\setlength{\tabcolsep}{3pt}
\renewcommand{\arraystretch}{1.12}
\caption{Decision Matrix for Problem Scope}
\label{tab:formal-problem-and-hypotheses-fp-06}
\begin{tabular}{|p{0.29\textwidth}|p{0.29\textwidth}|p{0.29\textwidth}|}
\hline
Decision & Criterion & Consequence for Next Stage \\\hline
Accept AANEC as matter assumption & The integral relation is stated and kept distinct from pointwise NEC/SEC. & Definitions and forward derivation may use it as an averaged energy premise. \\\hline
Require independent focusing assumption & AANEC alone is not treated as implying trapped-surface formation. & The proof chain must carry A5 as an explicit geometric hypothesis. \\\hline
Exclude astrophysical mass-gap inference & The observed 2-5 solar mass gap is not assumed to follow directly from the theorem. & Later sections must separate formal compactness bounds from stellar-evolution explanations. \\\hline
\end{tabular}
\end{table*}~\cite{cite_72d7ff391b29c495,cite_7a407c865faafd11,cite_7a44e6730c44cb39,cite_cbbced417410bc6c,cite_f634820e22d7a238,cite_fd0f77b2a173c0ab}

\section{Study Design and Methods}
\subsection{Protocol Scope and Independent Focusing}
The unit of analysis is a formal proof and counterexample search over a restricted class of spherically symmetric stellar collapse models. The protocol replaces pointwise energy assumptions with the averaged null energy condition and an independent averaged null-convergence/focusing hypothesis, rather than treating the latter as a consequence of the former. This separation is motivated by the recognition that quantum fields routinely violate pointwise energy conditions, making averaged statements necessary for broader validity. The design therefore isolates the additional geometric assumption required to force trapped-surface formation and evaluates whether it can be closed without reintroducing the pointwise restrictions it was intended to relax.~\cite{cite_0ed84243cf171bc4}

\subsection{Averaged Focusing Dependency}
∫\_γ T\_ab k\textasciicircum{}a k\textasciicircum{}b dλ ≥ 0~\cite{cite_0ed84243cf171bc4}

\subsection{Operational Role of the Focusing Integral}
The displayed relation records the averaged null energy condition along a null generator with tangent vector k\textasciicircum{}a and affine parameter λ. The protocol requires a second, strictly stronger averaged focusing condition on the expansion scalar along the same generator; its precise mathematical form is not yet fixed and must be supplied by human review before any candidate witness can be adjudicated. Until that dependency is resolved, the integral above serves only as the baseline energy hypothesis, not as the operative focusing premise.

\begin{table*}[htbp]
\centering
\footnotesize
\setlength{\tabcolsep}{3pt}
\renewcommand{\arraystretch}{1.12}
\caption{Boundary Checks and Counterexample Screening}
\label{tab:study-design-and-methods-sd-4}
\begin{tabular}{|p{0.29\textwidth}|p{0.29\textwidth}|p{0.29\textwidth}|}
\hline
Check & Admissible if & Rejected if \\\hline
Spherical symmetry & Metric and pressure are isotropic & Rotation, magnetic fields, or nonspherical dynamics dominate \\\hline
Null generator completeness & Generators extend to infinite affine parameter & Incompleteness or quantum stress-energy violations obstruct the integral \\\hline
Averaged focusing & The independent focusing hypothesis holds & Only the baseline energy condition holds \\\hline
Stable remnant predicate & Object is a regular equilibrium below the compactness limit & A witness satisfies all premises yet evades the trapped-surface conclusion \\\hline
\end{tabular}
\end{table*}

\subsection{Comparators, Robustness, and Transition}
\paragraph{Conditional outcome branch.} The design compares the averaged-focusing route against the classical pointwise assumption, treating the observed 2-5 solar mass gap as an astrophysical consistency constraint rather than a theorem-derived prediction. Explosion energetics, fallback, equation-of-state physics, pair-instability processes, and binary evolution are declared confounders that can populate or erase the gap independently of any compactness bound. Expected outcomes remain unobserved: the mechanism-support branch would arise only if the prespecified analysis is consistent with the declared relation while planned controls do not favor a stated alternative; a null or contradictory branch would indicate the relation is unsupported within the design boundary; and an uninformative branch would follow a validity or data-quality failure. Resolution of the open measurement and reproducibility decisions is required before the forward derivation can proceed to a decision rule.~\cite{cite_72d7ff391b29c495,cite_7a407c865faafd11}

\section{Expected Outcome Branches and Conditional Conclusions}
\subsection{Decision Criterion for the Formal Branch}
This section converts the scoped premises inherited from the forward derivation into a prespecified decision criterion for the proposed formal program. The derivation supplies the chain from averaged null-convergence and focusing (S1) through trapped-surface identification (S2) to the stable-remnant exclusion bound (S3-S4), while the counterexample search supplies the boundary witnesses that the criterion must reject or accommodate. The unique contribution here is to fix what counts as a supportive formal outcome before any proof or disproof is written: a branch is supportive only when the derivation closes for the declared admissible metric class and no witness survives all six assumptions. This keeps the expected outcomes separate from the astrophysical mass-gap application, which remains a consistency constraint rather than a theorem-derived prediction.

\begin{table*}[htbp]
\centering
\footnotesize
\setlength{\tabcolsep}{3pt}
\renewcommand{\arraystretch}{1.12}
\caption{Conditional Outcome Decision Matrix}
\label{tab:expected-outcomes-eo-02}
\begin{tabular}{|p{0.23\textwidth}|p{0.23\textwidth}|p{0.23\textwidth}|p{0.23\textwidth}|}
\hline
Branch & Prespecified Trigger & Interpretation & Decision Consequence and Next Action \\\hline
supports\_mechanism & The derivation closes for every admissible metric in the restricted class, and proposed witnesses fail at least one assumption. & The scoped relation is supported within the declared boundary; support is not proof of general validity. & Carry the qualified bound forward; test the most consequential boundary condition and replicate under independently confirmed formal assumptions. \\\hline
partial\_or\_heterogeneous & The bound holds for some admissible subclasses but fails or weakens for others, such as differing regularity or completeness regimes. & The relation is conditional; no universal exclusion is warranted. & Predefine and check plausible moderators, then improve coverage and comparability across the declared metric subclasses. \\\hline
null\_or\_contradictory & A witness satisfies all stated assumptions yet evades trapped-surface formation, or the comparison favors a declared alternative. & The proposed relation is not supported in this design boundary; absence of support is not general proof of absence. & Audit construct validity and comparison adequacy, then revise the mechanism or boundary claim before another design iteration. \\\hline
uninformative\_or\_invalid & Missing formal definitions prevent evaluation of the criterion, or a quality-control or validity failure blocks interpretation. & No scientific conclusion is warranted from the planned formal design. & Resolve the identified validity or definitional failure and obtain human confirmation of the prerequisites before repeating the design. \\\hline
\end{tabular}
\end{table*}

\subsection{Formal Threshold Relation}
C\_\{\textbackslash{}mathrm\{crit\}\} = \textbackslash{}inf \textbackslash{}\{ C \textbackslash{}in \textbackslash{}mathcal\{C\}\_\{\textbackslash{}mathrm\{adm\}\} : \textbackslash{}exists\textbackslash{}, \textbackslash{}mathcal\{S\} \textbackslash{}text\{ a trapped surface under \} A1\textbackslash{}text\{-\}A6 \textbackslash{}\}

\subsection{Reading the Threshold and the Supportive Branch}
The proposed threshold above is a decision object, not a computed number: it is the infimum of the compactness values in the admissible class for which a trapped surface is forced under the six assumptions. The supportive branch is triggered only when this infimum is attained within the declared model class and the derivation S1-S4 closes without a surviving witness. The counterexample analysis is decisive here because both proposed witnesses are currently assessed as failing at least one assumption, so they do not yet contradict the bound; they instead locate the boundary where the criterion must be tested. The averaged-focusing premise is essential because AANEC alone may not yield the integrated convergence that the derivation requires, a point consistent with the survey's finding that quantum fields systematically violate pointwise conditions and force reliance on averaged statements. The exclusion is therefore conditional on the geometric and stability assumptions, and it does not by itself reproduce the observed 2-5 solar mass gap, whose location is governed by near-final core mass and explosion physics rather than by the theorem alone.~\cite{cite_0ed84243cf171bc4,cite_7a407c865faafd11}

\subsection{Competing-Explanation and Invalidity Branches}
The heterogeneous and null branches are read against the declared alternative explanations. If the bound varies across regularity or completeness regimes, the outcome is heterogeneous and no universal exclusion follows; if a witness satisfies all assumptions while evading the trapped surface, the relation is not supported in this boundary. Either way, the astrophysical mass gap remains non-discriminating on its own: explosion energetics, fallback, equation-of-state physics, rotation, magnetic fields, nonspherical dynamics, and binary evolution can populate or reshape the observed distribution independently of a compactness theorem. The invalid branch is the most consequential for this route, because the threshold object cannot be evaluated while the averaged-focusing premise and the compactness definition remain unspecified; a formal design that cannot state its own criterion yields no interpretable evidence. The next action is therefore not to assert the bound but to resolve the definitional prerequisites, after which the decision matrix can be applied to a concrete proof or counterexample attempt.~\cite{cite_72d7ff391b29c495,cite_c369df652181e8f7}

\subsection{Transition to Review}
This conditional-outcome structure hands the argument to the review stage with a concrete dependency: the decision matrix is executable only once the formal prerequisites it invokes are fixed. The review section owns the item-by-item release criteria and human-confirmation requirements, so this route states the consequence rather than the ledger. The transition is that the supportive, heterogeneous, and null branches each presuppose a well-defined threshold and a well-defined averaged-focusing condition; until those are supplied, the invalid branch dominates and no formal conclusion can be released.

\section{Risks, Limitations, and Human Review Requirements}
The proposed AANEC-relaxed trapped-surface theorem is confined to classical general relativity under spherical symmetry, isotropic pressure, global hyperbolicity, and complete null generators. It does not quantitatively derive the observed 2-5 solar mass gap, nor does it universally exclude all non-black-hole remnants. The theorem must be restricted if quantum stress-energy effects violate averaged conditions, or if nonspherical, rotating, or magnetic dynamics dominate collapse outcomes. The addition of an explicit averaged null-convergence condition distinguishes this framework from standard Penrose proofs, but it may be as restrictive as the original pointwise assumptions in some regimes. Alternative stellar-physics mechanisms, including explosion energetics, fallback, equation-of-state uncertainties, and binary evolution, remain confounders for the mass-gap inference.~\cite{cite_0ed84243cf171bc4}

Two unverified candidate counterexamples challenge the exclusion of stable non-black-hole remnants. CE1 proposes a hypothetical static, spherically symmetric regular metric with high compactness that satisfies AANEC but violates the explicit averaged focusing condition. CE2 suggests regular black hole models, such as Hayward or Bardeen metrics, where trapped surfaces form but singularities are resolved, challenging the 'non-black-hole' classification. Both candidates currently fail to satisfy all theorem assumptions, particularly regarding global hyperbolicity and the independent averaged focusing constraint. Because the mathematical form of the averaged focusing condition is undefined, these witnesses cannot be rigorously evaluated. Formal verification requires resolving these structural ambiguities before any counterexample can be deemed valid or invalid.

\begin{table*}[htbp]
\centering
\footnotesize
\setlength{\tabcolsep}{3pt}
\renewcommand{\arraystretch}{1.12}
\caption{Decision matrix}
\label{tab:risk-limitations-and-review-decision-matrix}
\begin{tabular}{|p{0.29\textwidth}|p{0.29\textwidth}|p{0.29\textwidth}|}
\hline
Decision Point & Status & Required Action \\\hline
Formal Reasoning Plan & Review Required & Qualified human review of the final canonical plan status \\\hline
Misner-Sharp Compactness & Needs Input & Define formal equation for C and threshold C\_max \\\hline
Averaged Focusing (A5) & Needs Input & Specify mathematical form distinguishing A5 from AANEC \\\hline
Remnant Stability (A6) & Needs Input & Define stability predicate to exclude regular black holes \\\hline
Counterexample Validity & Needs Input & Determine if CE1 and CE2 genuinely fail A5 or if A5 is ill-defined \\\hline
\end{tabular}
\end{table*}

\paragraph{Human-review checklist.} 1. Qualified human review of the formal reasoning plan must be completed to approve the overall logical structure. 2. The exact mathematical formulation of the explicit averaged null-convergence condition must be provided to resolve proof obligations and evaluate candidate witnesses. 3. The formal definition of the Misner-Sharp compactness parameter must be established to calculate precise mass and compactness thresholds. 4. The stability predicate for non-black-hole remnants must be formally defined to distinguish stable equilibrium objects from regular black hole models like Hayward or Bardeen metrics.

\section{Definitions, Propositions, and Proof Obligations}
\subsection{Scope of the Formal Ledger}
This section establishes the definitional and propositional backbone for the proposed AANEC-relaxed trapped-surface theorem. It translates the scoped premises received from the formal problem stage into a rigorous symbol domain, candidate definitions, and unresolved proof obligations. The intellectual labor here is to fix the mathematical conventions—specifically the compactness and focusing parameters—so that the subsequent forward-derivation stage can operate on a stable formal object rather than an ambiguous verbal claim. No empirical results are claimed; all constructs are proposed formalizations subject to human verification.

\subsection{Primary Definition and Symbol Ledger}
\paragraph{Definition.} The following ledger assigns canonical symbols to the physical and geometric quantities. Items marked with a dependency flag represent incomplete formalizations where the mathematical form is not yet fixed in the input records; these are not invented here but identified as required human inputs.

- **M (D1)**: Mass of the collapsing matter or remnant object. Status: Candidate formalization. - **r (D2)**: Areal radius or radial coordinate. Status: Candidate formalization. - **theta (D3)**: Expansion scalar of a null geodesic congruence. Status: Candidate formalization. - **C (D4)**: Misner-Sharp compactness parameter. Status: **Needs Human Input**. The specific formula for effective mass within radius r in the chosen metric is missing, preventing precise threshold calculation. - **T (D5)**: Stress-energy tensor. Status: Candidate formalization. - **g (D6)**: Spacetime metric. Status: Candidate formalization. - **k (D7)**: Tangent vector to null geodesic generators. Status: Candidate formalization.

\subsection{Proposed Threshold Relation}
\begin{equation}
C(r) >= C_threshold(M, EOS)
\label{eq:definitions-and-propositions-def-equation}
\end{equation}

\subsection{Interpretation of the Threshold Relation}
The relation above represents the proposed compactness criterion. The left-hand side, C(r), is the Misner-Sharp compactness parameter at areal radius r. The right-hand side, C\_threshold, is a model-dependent bound determined by the mass M and the equation of state (EOS). This equation serves as the formal anchor for the trapped-surface identification step. It is critical to note that the explicit mathematical definition of C(r) remains unresolved (see D4 above), meaning this relation is currently a structural placeholder for a rigorous inequality that must be instantiated by human input. See Eq.~\eqref{eq:definitions-and-propositions-def-equation}.

\subsection{Assumption Ledger and Boundary Conditions}
\begin{itemize}
\item The theorem's validity is strictly confined to the following declared assumptions. The distinction between A2 and A5 is a critical unresolved dependency.
\item **A1**: Spacetime and matter fields satisfy classical General Relativity with spherically symmetric collapse and isotropic pressure.
\item **A2**: Matter fields satisfy the Averaged Null Energy Condition (AANEC) along complete null generators.
\item **A3**: Spacetime is globally hyperbolic and sufficiently regular for trapped-surface arguments.
\item **A4**: Null generators in the relevant region are complete.
\item **A5**: An explicit averaged null-convergence/focusing condition holds. **Status: Needs Human Input**. The mathematical form distinguishing this from A2 is missing.
\item **A6**: Non-black-hole remnants are modeled as stable equilibrium objects obeying a bounded compactness or TOV-like limit.
\item *Failure Condition**: The theorem fails if AANEC is violated, if A5 fails, if null generators are incomplete, or if quantum stress-energy effects prevent pointwise focusing.
\end{itemize}

\subsection{Candidate Propositions}
\paragraph{Proposition (proposed).} Two formal propositions are proposed to structure the argument. These are unverified conjectures within the restricted model class.

**P1 (Trapped-Surface Formation)**: In spherically symmetric stellar collapse satisfying AANEC, global hyperbolicity, regular matter fields, complete null generators, and an explicit averaged null-convergence/focusing condition, reaching a model-dependent compactness threshold forces trapped-surface formation.

**P2 (Remnant Exclusion)**: Stable non-black-hole remnants above a compactness-dependent mass scale are excluded within the restricted model class defined by AANEC and averaged focusing.

\subsection{Unresolved Proof Obligations}
\begin{itemize}
\item The following obligations must be discharged for P1 and P2 to become verified theorems. These are currently marked as unresolved due to missing formal inputs.
\item **PO1**: Establish the mathematical form of the averaged null-convergence condition (A5) and prove it implies negative expansion theta along null generators. This is blocked by the missing definition of A5.
\item **PO2**: Prove that negative expansion theta combined with the compactness threshold implies the existence of a trapped surface. This requires rigorous geometric analysis of the metric g.
\item **PO3**: Derive the exclusion bound for stable non-black-hole remnants. This requires linking the trapped-surface result to the stability predicate A6 and the compactness definition D4.
\end{itemize}

\begin{table}[htbp]
\centering
\small


\caption{Decision Matrix for Formal Inputs}
\label{tab:definitions-and-propositions-def-table}
\begin{tabular}{|p{0.29\linewidth}|p{0.29\linewidth}|p{0.29\linewidth}|}
\hline
Decision Point & Condition & Action \\\hline
Compactness Definition & Formula for C(r) missing & **Needs Human Input** (D4) \\\hline
Focusing Condition & Form of A5 distinct from A2 missing & **Needs Human Input** (PO1) \\\hline
Counterexample Search & Witness fails A5 & **Invalid Counterexample** \\\hline
Counterexample Search & Witness satisfies all A1-A6 & **Rejects Theorem** \\\hline
\end{tabular}
\end{table}

\section{Forward Derivation and Counterexample Search Plan}
\subsection{Derivation Architecture and Scope}
This section advances the argument from the scoped premises supplied by the definitions ledger into a proposed formal chain. The unique contribution is the operationalization of the Raychaudhuri-type focusing mechanism coupled with the Misner-Sharp compactness criterion, distinguishing admissible counterexamples from failures of the explicit averaged null-convergence assumption. The derivation proceeds through four numbered steps, culminating in the exclusion bound for stable non-black-hole remnants. Crucially, the explicit averaged focusing condition is treated as an independent geometric hypothesis rather than a consequence of AANEC alone, a distinction motivated by the systematic violation of pointwise energy conditions in quantum field theory contexts (W3009127927).~\cite{cite_0ed84243cf171bc4}

\subsection{Consumer Premise Subset}
\begin{itemize}
\item The forward derivation consumes the following specific premises from the definition ledger to construct the lemma chain:
\item A2: Matter fields satisfy AANEC along complete null generators.
\item A5: An explicit averaged null-convergence/focusing condition holds along relevant generators.
\item A6: Non-black-hole remnants are modeled as stable equilibrium objects obeying a bounded compactness limit.
\item D4: Misner-Sharp compactness parameter (formal definition pending human input).
\item D7: Tangent vector to null geodesic generators.
\end{itemize}

\subsection{Integrated Focusing Relation}
\begin{equation}
dtheta/dlambda = -1/2 theta^2 - sigma_ab sigma^ab + omega_ab omega^ab - R_ab k^a k^b
\label{eq:forward-derivation-and-counterexamples-fd-equation-block}
\end{equation}

\subsection{Explanation of Derivation Step S1}
The proposed Raychaudhuri equation governs the evolution of the expansion scalar theta along null generators. Under the explicit averaged null-convergence condition (A5), the integrated effect of the Ricci term R\_ab k\textasciicircum{}a k\textasciicircum{}b drives theta negative within a finite affine parameter, provided shear dominates vorticity. This constitutes the proposed derivation step S1, which is unverified pending the closure of proof obligation PO1 regarding the precise mathematical form of A5. See Eq.~\eqref{eq:forward-derivation-and-counterexamples-fd-equation-block}.

\subsection{Proposed Exclusion Bound}
\paragraph{Proposition (proposed).} Proposition P2: Stable non-black-hole remnants above a compactness-dependent mass scale are excluded within the restricted model class defined by AANEC and averaged focusing.

\begin{table*}[htbp]
\centering
\footnotesize
\setlength{\tabcolsep}{3pt}
\renewcommand{\arraystretch}{1.12}
\caption{Counterexample Decision Matrix}
\label{tab:forward-derivation-and-counterexamples-fd-counterexample-matrix}
\begin{tabular}{|p{0.18\textwidth}|p{0.18\textwidth}|p{0.18\textwidth}|p{0.18\textwidth}|p{0.18\textwidth}|}
\hline
Witness ID & AANEC (A2) & Explicit Focusing (A5) & Stability (A6) & Decision \\\hline
CE1 (Hypothetical) & Satisfied & Violated & Boundary Conflict & Out-of-Domain \\\hline
CE2 (Regular BH) & Approximate & Unknown & Failed (BH not Remnant) & Invalid Counterexample \\\hline
\end{tabular}
\end{table*}

\subsection{Failure Conditions and Limitations}
The derivation fails or must be restricted if the explicit averaged null-convergence condition (A5) is proven to be mathematically equivalent to AANEC in the chosen metric class, rendering it redundant, or if it is too strong to be physically motivated. Furthermore, the exclusion bound is invalid if candidate witnesses (such as Hayward or Bardeen metrics) can be shown to satisfy A5 while remaining regular and stable non-black-hole remnants. Currently, the status of CE2 is marked as assumptions\_not\_satisfied because regular black holes typically contain horizons, thereby failing the definition of a non-black-hole remnant in A6.

\appendices
\section{Idea Source Checkpoints and Direction Selection Audit}
\subsection{Audit of Directional Snapshots}
The selected mechanism\_replacement direction is supported by three available audit snapshots that share a common intellectual core rather than a proven temporal sequence. Because chronological metadata is absent, these checkpoints function as parallel records of the proposal’s scope. Each snapshot converges on replacing pointwise energy assumptions with an Averaged Null Energy Condition (AANEC) supplemented by an independent averaged focusing hypothesis. This shared foundation forces a trapped surface when a model-dependent compactness threshold is exceeded, thereby excluding stable non-black-hole remnants above a specific mass scale. The consistent presence of this mechanism across the candidate, portfolio, and direction files confirms that the core theoretical argument was stable at the time of selection, even if the formal derivation remains incomplete.~\cite{cite_0ed84243cf171bc4}

\subsection{Retained Scope Corrections}
A critical decision retained across all snapshots is the demotion of the observed 2-5 solar mass gap from a direct theorem prediction to an astrophysical consistency constraint. The proposal explicitly acknowledges that the gap may be governed by supernova explosion physics, fallback, equation-of-state uncertainties, rotation, magnetic fields, nonspherical dynamics, or binary interactions. This scope correction is essential for maintaining the separation between formal claims and empirical observations. By treating the mass gap as a confounder rather than a necessary consequence of the AANEC theorem, the design prevents the formal model from being falsified by astrophysical complexities that lie outside its restricted spherical, non-rotating domain. The snapshots uniformly preserve this boundary, ensuring the theorem is evaluated on its geometric merits rather than its ability to reproduce population demographics.~\cite{cite_7a407c865faafd11,cite_f634820e22d7a238}

\subsection{Preserved Open Questions}
The audit reveals that specific mathematical gaps were preserved rather than resolved during direction selection. Most notably, the explicit averaged null-convergence condition and the formal definition of the Misner-Sharp compactness parameter remain undefined in the upstream records. These omissions are not treated as errors to be silently patched but as open design questions requiring human input. The snapshots consistently flag these dependencies, indicating that the proposal’s viability hinges on closing the gap between the declared AANEC assumption and the additional focusing condition required for the Raychaudhuri-type derivation. Preserving these unknowns ensures that subsequent formal work addresses the actual missing links rather than assuming a completed proof.

\begin{table*}[htbp]
\centering
\footnotesize
\setlength{\tabcolsep}{3pt}
\renewcommand{\arraystretch}{1.12}
\caption{Convergence of Audit Snapshots}
\label{tab:appendix-idea-evolution-checkpoint-convergence-table}
\begin{tabular}{|p{0.29\textwidth}|p{0.29\textwidth}|p{0.29\textwidth}|}
\hline
Decision Dimension & Status Across Snapshots & Consequence for Formal Work \\\hline
Mass Gap Interpretation & Consistent: Astrophysical constraint only & Empirical falsifiers must target the theorem, not the population distribution \\\hline
AANEC + Focusing Mechanism & Consistent: Core replacement strategy & Raychaudhuri derivation must explicitly separate AANEC from added focusing \\\hline
Misner-Sharp Definition & Consistent: Undefined & Threshold calculations are blocked until human input specifies the metric \\\hline
Counterexample Validity & Consistent: Unverified & Candidate witnesses cannot be dismissed without formal A5 definition \\\hline
\end{tabular}
\end{table*}~\cite{cite_0ed84243cf171bc4}

\section{Variables, Symbols, and Operational Definitions}
\subsection{Core Geometric and Matter Variables}
\paragraph{Definition.} The formal argument operates on a restricted set of geometric and matter variables. The spacetime metric g and stress-energy tensor T define the background and source, while the null tangent vector k and expansion scalar theta characterize the congruence whose focusing is under investigation. The mass M and areal radius r provide the dimensional scales for the collapse. These variables are treated as candidate formalizations, requiring rigorous symbol-domain conventions to ensure that the proposed Raychaudhuri-type focusing chain remains mathematically well-posed. See Eq.~\eqref{eq:appendix-variables-and-definitions-block-2}.

\subsection{Averaged Null Energy Condition (AANEC)}
\begin{equation}
int_{-infty}^{infty} T_{ab} k^a k^b lambda d lambda >= 0
\label{eq:appendix-variables-and-definitions-block-2}
\end{equation}

\subsection{Operational Dependencies and Unresolved Specifications}
While the AANEC is formally defined, its operational application requires specifying the equation of state (EOS) and the precise form of the explicit averaged null-convergence condition. The EOS models (V10) and the compactness threshold formula (V6) remain undefined, preventing the calculation of a numerical mass scale. Furthermore, the operational definitions for rotation (V11) and magnetic fields (V12) are absent, restricting the current formalism to spherically symmetric, non-rotating, and unmagnetized collapse. These dependencies must be resolved to transition from a qualitative theorem to a quantitative prediction.

\begin{table*}[htbp]
\centering
\footnotesize
\setlength{\tabcolsep}{3pt}
\renewcommand{\arraystretch}{1.12}
\caption{Variable Operationalization Status}
\label{tab:appendix-variables-and-definitions-block-4}
\begin{tabular}{|p{0.23\textwidth}|p{0.23\textwidth}|p{0.23\textwidth}|p{0.23\textwidth}|}
\hline
Variable & Role & Operational Status & Decision \\\hline
V6 (Compactness C) & Independent & Formula missing & Define Misner-Sharp C \\\hline
V10 (EOS) & Confounder & Models unspecified & Select nuclear physics parameters \\\hline
V11 (Rotation) & Confounder & Initial conditions undefined & Specify angular momentum profile \\\hline
V12 (Magnetic Field) & Confounder & Measurement undefined & Define field topology \\\hline
V7 (Remnant Mass) & Dependent & Measurement method unspecified & Choose radial velocity vs timing \\\hline
V8 (Explosion Energy) & Confounder & Metric (kinetic/total) undefined & Define energy budget \\\hline
\end{tabular}
\end{table*}

\section{Evidence Coverage, Unknown Items, and Review Checklist}
\subsection{Bounded Evidence Coverage}
The proposal's formal focus is anchored by the observation that pointwise energy conditions are systematically violated by quantum fields, making averaged statements such as the AANEC the appropriate classical limit (W3009127927). This motivates the theorem's replacement of pointwise NEC/SEC with AANEC plus an explicit averaged null-convergence condition. On the astrophysical side, the observed 2-5 solar mass gap is treated strictly as a consistency constraint, not a theorem-derived prediction. Population analyses of GWTC-3 (W4226083230) and the anomalous 2.6 solar mass secondary in GW190814 (W3036761016) demonstrate that the gap's empirical status remains contested by selection effects and exotic formation channels. Similarly, the pair-instability gap inferred from GW190521 (W3081939279) and the lower-edge predictions of Farmer et al. (W2994776863) show that explosion energetics and fallback can populate or reshape the gap independently of any compactness bound. These sources justify the boundary condition that the mass-gap inference fails if stellar-physics mechanisms dominate, while leaving the formal theorem's scope intact.~\cite{cite_0ed84243cf171bc4,cite_f634820e22d7a238,cite_7a407c865faafd11,cite_cbbced417410bc6c,cite_fd0f77b2a173c0ab}

\begin{table*}[htbp]
\centering
\footnotesize
\setlength{\tabcolsep}{3pt}
\renewcommand{\arraystretch}{1.12}
\caption{Unknown-Item Consequences for the Formal Chain}
\label{tab:appendix-evidence-and-review-unknown-consequences}
\begin{tabular}{|p{0.29\textwidth}|p{0.29\textwidth}|p{0.29\textwidth}|}
\hline
Missing Input & Consequence for This Route & Required Decision \\\hline
Misner-Sharp compactness C (D4) & Prevents precise calculation of the threshold C\_max and the exclusion bound in S4. & Supply the effective-mass formula for the chosen metric class before deriving P2. \\\hline
Averaged focusing condition A5 & Renders the counterexample witness CE1 unverifiable, since its A5 check is false by default. & Specify the integral form of A5 to distinguish it from AANEC (A2). \\\hline
Stability predicate A6 & Creates a definitional conflict: an object above C\_max may fail 'IsStableCompactObject' rather than falsify the theorem. & Define the TOV-like bound independently of C\_max to avoid circular exclusion. \\\hline
Numerical verification scope & Leaves open whether Raychaudhuri focusing should be tested on discrete test metrics. & Decide if numerical checks supplement or replace the analytic proof obligation PO1. \\\hline
\end{tabular}
\end{table*}

\subsection{Release Criteria for Future Claims}
\paragraph{Human-review checklist.} The formal reasoning plan currently carries a status of requires\_human\_review (review-1). No proposition may advance beyond candidate\_formalization until the following gates are cleared. First, P1 and P2 must be labeled proposed and unverified; the terms proved, established, verified, demonstrated, or observed are forbidden until the proof obligations PO1-PO3 are closed. Second, the counterexample status must remain candidate\_found\_unverified: CE1 and CE2 fail assumption checks (A5 or A6) and therefore do not constitute valid counterexamples unless human review resolves the missing definitions. Third, the observed 2-5 solar mass gap must never be presented as direct evidence for the theorem; it is a design assumption subject to alternative explanations including rotation, magnetic fields, and binary evolution. Fourth, the outcome branch supports\_mechanism may only be triggered by a prespecified analysis consistent with the declared relation within the design boundary, and never by astrophysical observations alone. Release requires that all unknown items above be converted into formal definitions or explicitly scoped out as limitations.\section*{References}
\bibliographystyle{IEEEtran}
\bibliography{references}
\end{document}
